\documentclass{article}
\usepackage[top=2cm,bottom=2cm,left=2cm,right=2cm]{geometry}
\usepackage[dvipsnames,svgnames,x11names]{xcolor}
\usepackage{amsmath,amssymb,mathtools}
\usepackage{graphicx}
\usepackage{float}
\usepackage{cancel}
\usepackage{tcolorbox}
\usepackage{hyperref}
\usepackage{amsopn,amssymb}
\hypersetup{pdfauthor={Author}, pdftitle={Assignement 1}}

\setlength{\parindent}{0pt}
\setlength{\parskip}{6pt}

\begin{document}

% --- PAGE 1 ---
\begin{center}
    {\LARGE \underline{Assignement 1.}}
\end{center}

\vspace{0.5cm}

{\Large \underline{Classical Thermodynamics}}

\vspace{0.3cm}

\begin{enumerate}
    \item[\textbf{1.}] The basic 3 assumptions are:
    \begin{enumerate}
        \item[\textbf{1.}] monolayer coverage. Adsorption forms only one layer of molecules on the surface adsorption cannot occur on a site that is already occupied.
        \item[\textbf{2.}] the surface is uniform and homogenous, all binding sites have the same affinity and are energetically equivalent.
        \item[\textbf{3.}] No interaction between the molecules.
    \end{enumerate} 

    \item[\textbf{2.}] 
    \begin{equation*}
        q = \frac{q_{\text{sat}} \, K(T) \, P}{1 + K(T) \, P}
    \end{equation*}

    \begin{enumerate}
        \item[\textbf{a)}] $\displaystyle \lim_{P \to \infty} q = \lim_{P \to \infty} \frac{q_{\text{sat}} \, K(T) \, P}{1 + K(T) \, P} \longrightarrow \text{we notice the form } \lim_{x \to \infty} \frac{x}{1+x} = 1$
        
        \vspace{0.2cm}
        $\Rightarrow$ \fbox{\rule[-0.2cm]{0cm}{0.6cm} $\displaystyle \lim_{P \to \infty} q = q_{\text{sat}}$}

        \vspace{0.3cm}
        \item[\textbf{b)}] $\displaystyle \lim_{P \to 0} q = \lim_{P \to 0} \frac{q_{\text{sat}} \, K(T) \, P}{1 + K(T) \, P} \quad \text{at low } P \quad 1 + K(T) \, P \approx 1$

        \vspace{0.2cm}
        $\Rightarrow \text{at low pressure} \quad$ \fbox{\rule[-0.2cm]{0cm}{0.6cm} $q = q_{\text{sat}} \, K(T) \cdot P$}
    \end{enumerate}

    \vspace{0.3cm}
    \item[\textbf{3.}] At low pressure \quad $\displaystyle q = \underbrace{q_{\text{sat}} \cdot K(T)}_{K_H \text{ (Henry's coeff.)}} \cdot P$

    \vspace{0.3cm}
    \begin{equation*}
        K(T) = k_0 \exp\left(-\frac{\Delta H}{RT}\right) \Rightarrow \text{\fbox{\rule[-0.3cm]{0cm}{0.8cm} $\displaystyle K_H(T) = q_{\text{sat}} \, k_0 \exp\left(-\frac{\Delta H}{RT}\right)$}}
    \end{equation*}
\end{enumerate}

\newpage

% --- PAGE 2 ---
\begin{enumerate}
    \item[\textbf{4.}] \underline{\textbf{Step 1:}}
\end{enumerate}

\vspace{0.2cm}
\hspace{0.5cm} For each temp, a plot $q$ vs $P$ can be made to obtain $q_{\text{sat}}$, $k_{263}$, $k_{293}$, $k_{313}$

% \vspace{0.3cm}
% \begin{figure}[H]
%     \centering
%     \includegraphics[width=0.35\textwidth]{img-0.jpeg}
%     \caption{Plot of $q$ versus $P$}
% \end{figure}

\vspace{0.5cm}
\underline{\textbf{Step 2:}}

\vspace{0.3cm}
\begin{equation*}
    k(T) = k_0 \exp\left(-\frac{\Delta H}{RT}\right)
\end{equation*}

\begin{equation*}
    \ln k = \ln k_0 - \frac{\Delta H}{RT} \qquad \text{we can plot } \ln k \text{ vs } \frac{1}{T}
\end{equation*}

\vspace{0.3cm}
% \begin{minipage}[c]{0.45\textwidth}
%     \begin{figure}[H]
%         \centering
%         \includegraphics[width=\textwidth]{img-1.jpeg}
%         \caption{Plot of $\ln(K)$ versus $1/T$}
%     \end{figure}
% \end{minipage}
\hfill
\begin{minipage}[c]{0.5\textwidth}
    \begin{gather*}
        \text{getting the equation of the plot } y = mx + b \\
        \text{where } m = -\frac{\Delta H}{R} \quad \text{and} \quad b = \ln k_0 \\
        \text{thus determining } \Delta H \text{ and } k_0
    \end{gather*}
\end{minipage}

\newpage

% --- PAGE 3 ---
{\Large \underline{Statistical Thermodynamics:}}

\vspace{0.5cm}

\begin{enumerate}
    \item[\textbf{1.}] \textbf{Canonical ensemble}

    \vspace{0.3cm}
    \begin{tabular}{ll}
        It describes a system where & $N, V \text{ and } T \text{ are constant.}$
    \end{tabular}

    \vspace{0.2cm}
    The system can exchange energy with the environment surrounding it.

    \vspace{0.2cm}
    In the context of adsorption this means that each different state of the system has different energy.

    \vspace{0.5cm}
    \item[\textbf{2.}] $\displaystyle Q = \sum_v e^{-\beta E_v} \qquad \beta := \frac{1}{k_B T}$

    \vspace{0.5cm}
    \item[\textbf{3.}] $\displaystyle P_v = \frac{\exp(-\beta E_v)}{Q}$

    \vspace{0.5cm}
    \item[\textbf{4.}] 
    \begin{equation*}
        \langle S \rangle = k_B \ln Q + k_B T \left( \frac{\partial \ln Q}{\partial T} \right)_{N,V}
    \end{equation*}

    \begin{equation*}
        S = -k_B \sum_v P_v \ln P_v \quad \text{and} \quad P_v = \frac{e^{-\beta E_v}}{Q}
    \end{equation*}

    \begin{equation*}
        \Rightarrow S = -k_B \sum_v P_v (-\beta E_v - \ln Q) = k_B \beta \sum_v P_v E_v + k_B \ln Q \sum_v P_v \quad \text{knowing that } \sum_v P_v = 1 \text{ by definition}
    \end{equation*}

    \begin{equation*}
        \sum_v P_v E_v = \bar{E}_{\text{tot}}
    \end{equation*}

    \begin{align*}
        \Rightarrow S &= k_B \beta \bar{E} + k_B \ln Q \qquad \beta := \frac{1}{k_B T} \\
        &= k_B \ln Q + \frac{\bar{E}}{T} \qquad \text{using the expression given proofed later on: } \bar{E} = k_B T^2 \left( \frac{\partial \ln Q}{\partial T} \right)_{N,V}
    \end{align*}

    \vspace{0.3cm}
    \begin{equation*}
        \text{\fbox{\rule[-0.3cm]{0cm}{0.8cm} $\displaystyle \langle S \rangle = k_B \ln Q + k_B T \left( \frac{\partial \ln Q}{\partial T} \right)_{N,V}$}} \quad \blacksquare
    \end{equation*}
\end{enumerate}

% Continuation of handwritten physics notes (pages 4-5)
% \usepackage{cancel}
% \usepackage{tcolorbox}

% Page 4
\begin{flushleft}

\vspace*{0.2cm}

\begin{equation*}
P = k_B T \left( \frac{\partial \ln Q}{\partial V} \right)_{N,T}
\end{equation*}

\vspace{0.3cm}
Starting from:\quad $dE = T dS - P dV$

\begin{equation*}
P = T \left(\frac{\partial S}{\partial V}\right)_E - \left(\frac{\partial E}{\partial V}\right)_S
\end{equation*}

\vspace{0.2cm}
\begin{equation*}
\left.
\begin{aligned}
\frac{\partial S}{\partial V} &= k_B \left( \frac{\partial \ln Q}{\partial V} \right) + k_B T \left( \frac{\partial \ln Q}{\partial V} \right) \\[6pt]
\frac{\partial E}{\partial V} &= k_B T^2 \left( \frac{\partial \ln Q}{\partial V} \right)_{T,N}
\end{aligned}
\right\} \text{substituted in}
\end{equation*}

\vspace{0.2cm}
\begin{equation*}
\Rightarrow P = k_B T \left( \frac{\partial \ln Q}{\partial V} \right)_{N,T} + \cancel{k_B T^2 \left( \frac{\partial \ln Q}{\partial V} \right)_{N,T}} - \cancel{k_B T^2 \left( \frac{\partial \ln Q}{\partial V} \right)_{N,T}}
\end{equation*}

\vspace{0.2cm}
% \begin{equation*}
% \Rightarrow \quad \tcbox[colframe=red,colback=white,boxrule=1pt,arc=0mm,boxsep=3pt,left=6pt,right=6pt,top=4pt,bottom=4pt]{%
%   \displaystyle P = k_B T \left( \frac{\partial \ln Q}{\partial V} \right)%
% }   \quad \blacksquare
% \end{equation*}

% \vspace{0.6cm}

% \begin{equation*}
%     \mu = - k_B T \left( \frac{\partial \ln Q}{\partial N} \right)_{T,V}
% \end{equation*}

% \vspace{0.2cm}
% \begin{equation*}
% dE = T dS - P dV + \mu \, dN \quad \text{with } V \text{ constant}
% \end{equation*}

% \begin{equation*}
% \Rightarrow \mu = \frac{\partial E}{\partial N} - T \frac{\partial S}{\partial N}
% \end{equation*}

% \begin{equation*}
% \frac{\partial E}{\partial N} = k_B T^2 \left( \frac{\partial \ln Q}{\partial N} \right)_{T,V}
% \end{equation*}

% \begin{equation*}
% \frac{\partial S}{\partial N} = k_B \left( \frac{\partial \ln Q}{\partial N} \right)_{T,V} + k_B T \left( \frac{\partial \ln Q}{\partial N} \right)_{T,V}
% \end{equation*}

% \begin{equation*}
% \Rightarrow \mu = \cancel{k_B T^2 \left( \frac{\partial \ln Q}{\partial N} \right)_{T,V}} - k_B T \left( \frac{\partial \ln Q}{\partial N} \right)_{T,V} - \cancel{k_B T^2 \left( \frac{\partial \ln Q}{\partial N} \right)_{T,V}}
% \end{equation*}

% \vspace{0.2cm}
% \begin{equation*}
% \Rightarrow \quad \tcbox[colframe=red,colback=white,boxrule=1pt,arc=0mm,boxsep=3pt,left=6pt,right=6pt,top=4pt,bottom=4pt]{%
%   \displaystyle \mu = - k_B T \left( \frac{\partial \ln Q}{\partial N} \right)_{T,V}%
% \quad \blacksquare}
% \end{equation*}

\vspace{0.6cm}

\begin{equation*}
E = k_B T^2 \left( \frac{\partial \ln Q}{\partial T} \right)_{V,N}
\end{equation*}

\vspace{0.2cm}
\begin{equation*}
Q = \sum_\nu e^{-\beta E_\nu}
\end{equation*}

\begin{equation*}
\frac{\partial Q}{\partial \beta} = - \sum_\nu E_\nu e^{-\beta E_\nu} \quad \Rightarrow \quad \frac{\partial \ln Q}{\partial \beta} = \frac{1}{Q} \frac{\partial Q}{\partial \beta} \quad \Rightarrow \quad \frac{\partial \ln Q}{\partial \beta} = - \sum_\nu \frac{E_\nu e^{-\beta E_\nu}}{Q}
\end{equation*}

\end{flushleft}

\newpage

% Page 5
% \begin{flushleft}

\vspace*{0.2cm}

\begin{equation*}
\text{knowing } P_\nu = \frac{e^{-\beta E_\nu}}{Q} \quad \Rightarrow \quad \frac{\partial \ln Q}{\partial \beta} = - \sum_\nu P_\nu E_\nu \quad \text{since } E = \sum_\nu P_\nu E_\nu
\end{equation*}

\vspace{0.3cm}
\begin{equation*}
\Rightarrow \quad E = - \left( \frac{\partial \ln Q}{\partial \beta} \right)_{N,V}
\end{equation*}

\vspace{0.5cm}
\begin{equation*}
\text{with} \quad \beta = \frac{1}{k_B T} \quad \Rightarrow \quad \frac{\partial \beta}{\partial T} = - \frac{1}{k_B T^2}
\end{equation*}

\vspace{0.5cm}
\begin{equation*}
\text{using chain rule} \quad \frac{\partial \ln Q}{\partial T} = \frac{\partial \ln Q}{\partial \beta} \frac{\partial \beta}{\partial T} \quad \text{and substituting } \frac{\partial \beta}{\partial T} \text{ in}
\end{equation*}

\vspace{0.5cm}
\begin{equation*}
\Rightarrow \quad \frac{\partial \ln Q}{\partial T} = E \cdot \frac{1}{k_B T^2}
\end{equation*}

% \vspace{0.4cm}
% \begin{equation*}
% \Rightarrow \quad \tcbox[colframe=red,colback=white,boxrule=1pt,arc=0mm,boxsep=4pt,left=8pt,right=8pt,top=5pt,bottom=5pt]{%
%   \displaystyle E = k_B T^2 \left( \frac{\partial \ln Q}{\partial T} \right)_{V,N}%
% } \quad \blacksquare
% \end{equation*}

% \end{flushleft}

\end{document}